% Generated draft: requires mathematical and source review.
\documentclass[11pt]{article}
\usepackage[margin=25mm]{geometry}
\usepackage{fontspec}
\setmainfont{Latin Modern Roman}
\usepackage{amsmath,amssymb,mathtools,unicode-math}
\setmathfont{Latin Modern Math}
\usepackage{enumitem,needspace,xurl,hyperref}
\setlength{\parindent}{0pt}
\setlength{\parskip}{5pt}
\setlength{\emergencystretch}{3em}
\newcommand{\N}{\mathbb{N}}
\newcommand{\Z}{\mathbb{Z}}
\newcommand{\Q}{\mathbb{Q}}
\newcommand{\R}{\mathbb{R}}
\newcommand{\C}{\mathbb{C}}
\newcommand{\F}{\mathbb{F}}
\newcommand{\A}{\mathbb{A}}
\newcommand{\PP}{\mathbb{P}}
\newcommand{\E}{\mathbb{E}}
\newcommand{\eps}{\varepsilon}
\newcommand{\dd}{\,\mathrm{d}}
\begin{document}
\section*{UnsolvedMath ID: 800674}
\textbf{Source problem code:} catalogue entry
\subsection{Definitions and mathematical statement}
Let \(\mathbb D:=\{z\in\C:|z|<1\}\) and \(\mathbb D(r):=\{z\in\C:|z|<r\}\) for \(0<r<1\). The supplied excerpt proposes replacing a weighted current \(\tau_r^a\) by the unweighted current, called an \emph{Ahlfors current},
\[
 A(r):=\int_{\mathbb D(r)}\phi_a^*(\omega),
 \qquad
 \sigma_r^a:=\frac{1}{A(r)}(\phi_a)_*[\mathbb D(r)].
\]
Here \([\mathbb D(r)]\) denotes integration over the disc. The notation \(A(r)\) is retained from the source; its defining expression also involves \(a\). Whenever the maps, forms, and integrals are defined and \(0<A(r)<\infty\), the pushforward has the action
\[
 \langle\sigma_r^a,\eta\rangle
 =\frac{\displaystyle\int_{\mathbb D(r)}\phi_a^*(\eta)}
 {\displaystyle\int_{\mathbb D(r)}\phi_a^*(\omega)}
\]
on smooth compactly supported test forms \(\eta\) of bidegree \((1,1)\) on the ambient complex manifold.

\textbf{Question (Dinh--Sibony, Problem 1, 2018).} With the data and hypotheses intended in the source, does
\[
 \sigma_r^a\longrightarrow T\qquad\text{as }r\uparrow1
 \quad\text{for }T\text{-almost every }a?
\]
Under the usual interpretation as weak convergence of currents, this asks whether, for \(T\)-almost every \(a\), every such test form \(\eta\) satisfies
\[
 \lim_{r\uparrow1}
 \frac{\displaystyle\int_{\mathbb D(r)}\phi_a^*(\eta)}
 {\displaystyle\int_{\mathbb D(r)}\phi_a^*(\omega)}
 =\langle T,\eta\rangle.
\]
\textbf{Unresolved qualifications.} The supplied excerpt does not specify the ambient manifold, the foliation, the admissible points \(a\), the maps \(\phi_a\), the form \(\omega\), or the hypotheses and normalization imposed on \(T\). It also does not define the measure meant by \(T\)-almost every point. The positivity and finiteness needed for the displayed normalization are not stated. Thus this is a transcription and conditional interpretation of the question, not a fully specified conjecture with recoverable hypotheses. The limit is along all radii approaching \(1\), as written in the supplied question.
\subsection{Short English statement}
Dinh and Sibony ask whether the normalized unweighted integration currents of the discs \(\phi_a(\mathbb D(r))\) converge to \(T\) as \(r\uparrow1\), for \(T\)-almost every \(a\). The excerpt omits the surrounding geometric hypotheses and the meaning of this almost-everywhere condition. The registry labels the problem open; that label is supplied curation, not an independently verified status assessment.
\subsection{Sources}
\begin{itemize}
\item Supplied local UnsolvedMath registry record: permanent ID \(800674\), exact source code \texttt{catalogue entry}, title ``Holomorphic Foliation Theory --- Problem 1.'' The record attributes the problem to Dinh and Sibony, with proposed year 2018, and labels its status open.
\item Dinh and Sibony, \emph{Some Open Problems on Holomorphic Foliation Theory} (2018), Problem 1. Supplied source URL: \url{https://arxiv.org/abs/1807.05640}.
\item Supplied provenance records access on 2026-10-16 and extraction by \texttt{arxiv\_latex}. These are source-record provenance statements; the source was not independently accessed or verified for this draft.
\end{itemize}
\end{document}
